\subsection{Finiteness of the number of fixed sets}

\begin{Lemma}\label{finfix} Let $X$ be an affine scheme of finite type over ${\mathbb C}$ and
$G$ be a reductive group acting on $X$. Then up to the $G$-action there are finitely many subschemes
of $X$ of the form $X^g$, where $g\in G$ is a semisimple element, .
\end{Lemma}

\begin{proof} First consider the case when $G$ is diagonalizable (in which case every $g\in G$ is semisimple). In this case the lemma is standard (see, e.g., \cite[Proposition on p.~44]{Ste}). Namely, pick generators $f_1,\dots,f_n$ of ${\mathbb C}[X]$ such that $f_i(gx)=\chi_i(g)f_i(x)$ for some characters~$\chi_i$ of~$G$. Then $f_i(gx)-f_i(x)=(\chi_i(g)-1)f_i(x)$, so~$X^g$ is cut out by the equations $f_i(x)=0$ for those~$i$ for which $\chi_i(g)\ne 1$, which implies that there are finitely many subschemes~$X^g$.

The general case reduces to the diagonalizable case by using the following fact (see~\cite{Sie} and \cite[Section~1.14]{Lu}): for each connected component~$Z$ of~$G$ there is a diagonalizable subgroup $D_Z\subset G$ such that every semisimple element $g\in Z$ is $G^0$-conjugate to an element of~$D_Z$.\footnote{We are grateful to G.~Lusztig for this explanation.}
\end{proof}

\section{Classification of short star-products}\label{section3}

In this section we present the classification of nondegenerate short star-products, the idea of which was explained to the first author by Maxim Kontsevich.

\subsection{Twisted traces}

Consider the setting of Section~\ref{constru}. Let $g\colon {{\mathbf A}}\to {{\mathbf A}}$ be a filtration-preserving invertible linear map, and let $T\colon {{\mathbf A}}\to {\mathbb C}$ be a $g$-twisted trace, i.e., $T({\mathbf a}{\mathbf b})=T({\mathbf b}g({\mathbf a}))$ for all ${\mathbf a},{\mathbf b}\in {{\mathbf A}}$. Note that taking ${\mathbf b}=1$, we get $T(g({\mathbf a}))=T({\mathbf a})$, i.e., $T$ is $g$-invariant.

Define the bilinear form $(\,,\,)_T$ on ${\mathbf A}$ by the formula $({\mathbf a},{\mathbf b})_T:=T({\mathbf a}{\mathbf b})$.

\begin{Lemma}\label{ide} Let $K$ be the left kernel of the bilinear form $(\,,\,)_T$. Then $K$ is also the right kernel of $(\,,\,)_T$, and a $g$-invariant two-sided ideal in~${{\mathbf A}}$. In particular, if ${\mathbf A}$ is a simple algebra and $T\ne 0$ then $K=0$.
\end{Lemma}

\begin{proof} Let $K'$ be the right kernel of $(\,,\,)_T$. We have ${\mathbf b}\in K'$ if and only if $T({\textbf{ab}})=0$ for all ${\mathbf a}\in {{\mathbf A}}$. This is equivalent to saying that $T({\mathbf b}g({\mathbf a}))=0$ for all ${\mathbf a}\in {{\mathbf A}}$, which happens if and only if ${\mathbf b}\in K$. Thus, $K'=K$. Also ${\mathbf a}\in K$ if and only if $T({\mathbf b}g({\mathbf a}))=0$ for all ${\mathbf b}\in {\mathbf A}$, which happens if and only if $g({\mathbf a})\in K$, hence $K$ is $g$-invariant. Finally, if ${\mathbf a}\in K$ and ${\mathbf b}\in {{\mathbf A}}$
then $({\textbf{ab}},{\textbf{c}})_T= ({\textbf{a}},{\textbf{bc}})_T=0$ and $({\textbf{c}},{\textbf{ba}})_T=({\textbf{cb}},{\textbf{a}})_T=0$, so ${\textbf{ab}},{\textbf{ba}}\in K$, i.e., $K$ is a two-sided ideal, as desired.
\end{proof}

By Lemma~\ref{ide}, it makes sense to speak of the kernel of $(\,,\,)_T$ (without specifying if it is left or right).

\begin{Lemma} If the kernel of $(\,,\,)_T$ vanishes then $g$ is an algebra automorphism.
\end{Lemma}

\begin{proof}We have
\begin{gather*}
T({\mathbf c}g({\mathbf a}{\mathbf b}))=T({\textbf{abc}})=T({\textbf{bc}}g({\mathbf a}))=T({\mathbf c}g({\mathbf a})g({\mathbf b})),
\end{gather*}
hence $g({\mathbf{ab}})=g({\mathbf a})g({\mathbf b})$.
\end{proof}

\subsection{Nondegenerate short star-products and nondegenerate twisted traces}

We will say that $T$ is {\it nondegenerate} if the restriction of the bilinear form $({\mathbf a},{\mathbf b})_T:=T({\mathbf a}{\mathbf b})$ to~$F_i{{\mathbf A}}$ is nondegenerate for each~$i$. Note that this is a strictly stronger condition than vanishing of the kernel of $(\,,\,)_T$ (for example, it implies that $T(1)\ne 0$).

Suppose that $T$ is nondegenerate and $s$-invariant (in this case $g$ commutes with~$s$). Then the bilinear form $(\,,\,)_T$ defines a canonical orthogonal complement ${\mathbf A}^T_i$ to $F_{i-1}{{\mathbf A}}$ in $F_i{{\mathbf A}}$. Note that the right and left orthogonal complement coincide since $g$ is filtration-preserving. We have a natural isomorphism $\theta_i^T\colon {\mathbf A}_i^T\to A_i$. Define the quantization map $\phi^T\colon A\to {{\mathbf A}}$ by $\phi^T|_{A_i}:=\big(\theta_i^T\big)^{-1}$. Clearly, $\phi^T$ is $s$-invariant. Thus, $\phi^T$ defines a star-product $*$ on $A$. Clearly, this star-product does not depend on the normalization of~$T$, so we will always normalize~$T$ so that $T(1)=1$ (which is possible since for a nondegenerate~$T$, one has $T(1)\ne 0$).

\begin{Proposition}[M.~Kontsevich]\label{kon} The star-product $*$ is short and nondegenerate. Moreover, any nondegenerate short star-product on~$A$ is obtained in this way from a unique nondegenerate~$T$ such that $T(1)=1$. In other words, this correspondence is a bijection between short nondegenerate star-products on $A$ and filtered quantizations ${{\mathbf A}}$ of~$A$ equipped with a filtration-preserving automorphism~$g$ commuting with $s$ and a nondegenerate $s$-invariant $g$-twisted trace~$T$ such that $T(1)=1$.
\end{Proposition}

\begin{proof} We have a nondegenerate (in general, nonsymmetric) inner product on~$A$ given by $\beta(a,b):=(\phi(a),\phi(b))_T$, and the spaces $A_i$ are orthogonal under this inner product. Moreover, it is clear that $\beta(a*b,c)=\beta(a,b*c)$. Suppose $a\in A_m$, $b\in A_n$. Let us show that every homogeneous component of $a*b$ has degree $\ge |m-n|$. First assume that $m>n$. It suffices to show that for any homogeneous $c\in A$ of degree $d<m-n$, one has $\beta(a*b,c)=0$. But we have $\beta(a*b,c)=\beta(a,b*c)$, and $b*c$ does not have any terms of degree bigger than $n+d<m$.
Hence $\beta(a*b,c)=\beta(a,b*c)=0$, as desired. A similar argument applies if $m<n$: we have $\beta(c,a*b)=\beta(c*a,b)=0$. Thus, the star-product~$*$ is short. Moreover, we have
\begin{gather*}
\langle a,b\rangle={\rm CT}(a*b)=\beta(a*b,1)=\beta(a,b),
\end{gather*}
which implies that $\langle\,,\,\rangle$ (and hence~$*$) is nondegenerate.

Conversely, assume that $*$ is a nondegenerate short star-product on~$A$. Then the quantization~${{\mathbf A}}$ is identified with $(A,*)$ using the quantization map~$\phi$. Set $T({\mathbf a}):={\rm CT}\big(\phi^{-1}({\mathbf a})\big)$. Then the form $({\textbf{a}},{\textbf{b}})_T:=T({\textbf{ab}})$ is nondegenerate on each $F_i{{\mathbf A}}$. Thus there exists a unique linear automorphism~$g_i$ of $F_i{{\mathbf A}}$ such that $T({\textbf{ab}})=T({\mathbf b}g_i({\mathbf a}))$. Moreover, it is clear that~$g_i$ preserves~$F_{i-1}{{\mathbf A}}$, and $g_i|_{F_{i-1}{{\mathbf A}}}=g_{i-1}$. Thus, the maps~$g_i$ define a~filtration-preserving linear automorphism $g\colon {{\mathbf A}}\to {{\mathbf A}}$.

Finally, it is easy to check that these assignments are mutually inverse, as claimed.
\end{proof}

\begin{Remark} If $T$ is any $s$-invariant linear functional on ${\mathbf A}$ such that the form $({\mathbf a},{\mathbf b})_T=T({\mathbf a}{\mathbf b})$ is nondegenerate on each $F_i{\mathbf A}$ (which clearly happens for ``random'' $T$) then for each~$i$ we can define a linear automorphism $g_i$ of $F_i({\mathbf A})$ such that $T({\textbf{ab}})=T({\mathbf b}g_i({\mathbf a}))$. However, in general $g_i|_{F_{i-1}{\mathbf A}}\ne g_{i-1}$ and $g_i$ does not preserve~$F_{i-1}{\mathbf A}$ (so the automorphism~$g$ of ${\mathbf A}$ restricting to~$g_i$ on each $F_i{\mathbf A}$ is not defined). Thus, the left and right orthogonal complements of~$F_{i-1}{\mathbf A}$ in~$F_i{\mathbf A}$ do not coincide, and one cannot define a quantization map giving a short star-product. The condition that~$g$ is well defined is a very strong restriction of~$T$, and we will see that it usually leads to a finite dimensional space of possible~$T$.
\end{Remark}

Let ${\mathbf A}g$ be the ${\mathbf A}$-bimodule which is ${\mathbf A}$ as a left module and ${\mathbf A}$ with action twisted by $g$ as a right module.

\begin{Corollary}\label{bun} Let $S=S({\mathbf A})$ be the set of nondegenerate short star-products corresponding to a filtered quantization ${{\mathbf A}}$ of $A$. Then we have a map $\pi\colon S\to \operatorname{Aut}({{\mathbf A}})$ from $S$ to the group of filtration-preserving $s$-invariant automorphisms of ${{\mathbf A}}$ whose fiber $\pi^{-1}(g)$ is a subset of the space $(HH_0({{\mathbf A}},{{\mathbf A}}g)^s)^*$ dual to the $s$-invariants in the Hochschild homology $HH_0({{\mathbf A}},{{\mathbf A}}g)$.
\end{Corollary}

\begin{proof} This follows immediately from Proposition~\ref{kon}, since by definition, an $s$-invariant $g$-twisted trace on ${{\mathbf A}}$ is an element of $(HH_0({{\mathbf A}},{{\mathbf A}}g)^s)^*$.
\end{proof}

\begin{Corollary} \label{bun1} Suppose $A=\mathcal{O}(X)$, where $X$ is a conical symplectic
singularity. Then any nondegenerate short star-product corresponding to $g\in \operatorname{Aut}({{\mathbf A}})$
is invariant with respect to the connected component~$Z_g^0$ of the centralizer~$Z_g$
of~$g$ in $\operatorname{Aut}({{\mathbf A}})$. In particular, any nondegenerate even short star-product
is invariant under $\operatorname{Aut}({\mathbf A})^0$.
\end{Corollary}

\begin{proof} It suffices to show that the Lie algebra $\operatorname{Lie}Z_g$ acts trivially on
$HH_0({\mathbf A},{\mathbf A} g)$. Let $z\in \operatorname{Lie}Z_g$. By Lemma~\ref{inn}, there exists a unique ${\mathbf z}\in F_2{\mathbf A}$ up to adding a constant such that $z({\mathbf a})=[{\mathbf z},{\mathbf a}]$ for all ${\mathbf a}\in {\mathbf A}$. Then $g({\mathbf z})={\mathbf z}+C$ for some constant $C$, so $C=g({\mathbf z})-{\mathbf z}$. Thus if $C\ne 0$ then for any $g$-twisted trace $T$ we have $T(1)=0$, so there are no nondegenerate short star-products at all. On the other hand, if $C=0$ then we have $[{\mathbf z},{\mathbf a}]={\mathbf z}{\mathbf a}-{\mathbf a}g({\mathbf z})$, so ${\mathbf z}$ acts trivially on $HH_0({\mathbf A},{\mathbf A} g)$, as claimed. The last statement follows from the fact that in the even case $g=s$, and~$s$ by definition is central in~$\operatorname{Aut}({{\mathbf A}})$.
\end{proof}

